Se diseña la capa de datos de un copiloto RAG interno para una distribuidora
mayorista. La información está repartida entre documentos, procedimientos,
condiciones comerciales y datos operativos. Por eso, encontrar una respuesta
confiable puede llevar tiempo y depender del conocimiento informal de otras
personas.

Se usa PostgreSQL como motor principal, con las extensiones
\texttt{pgvector} y \texttt{btree\_gist}. La misma base guarda datos
operativos, metadatos documentales, fragmentos, embeddings, permisos y
auditoría. Así se aplican la vigencia y los permisos antes de buscar por
similitud, y las fuentes vencidas o no autorizadas quedan fuera del ranking.

La implementación incluye datos sintéticos reproducibles, siete consultas
funcionales y dos pruebas de seguridad y auditoría. Las restricciones, las
políticas RLS (seguridad a nivel de fila) y la evidencia de cada respuesta se
validan mediante dos cargas limpias equivalentes. La prueba termina en la capa
de datos: no incluye interfaz, API ni un modelo de lenguaje real.
